\documentclass[10pt,a4paper]{amsart}
\usepackage[margin=19mm]{geometry}
\usepackage{amsmath,amssymb,mathtools}
\usepackage[scr=rsfs,cal=euler]{mathalfa}
\usepackage{xfrac}
\usepackage[unicode,hidelinks,bookmarks=false]{hyperref}
\newcommand{\ie}{i.e.\ }
\newcommand{\cf}{cf.\ }
\newcommand{\resp}{respectively\ }
\newcommand{\Subsection}{\subsection}
\providecommand{\nobreakdash}{\nobreak\hbox{-}}
\setlength{\emergencystretch}{2em}
\multlinegap0pt
\usepackage{amssymb}
\usepackage{xcolor}
\usepackage{tikz}
\usepackage[shortlabels]{enumitem}
\newtheorem{lemma}{Lemma}
\def\Eig{{\rm Eig}}
\def\hyper{{\mathcal{H}}}
\title{Induced dynamics and quasifactors for minimal equicontinuous actions on Stone spaces}
\author{Mar\'ia Isabel Cortez, Till Hauser}
\date{Source: arXiv:2602.11034v1 (CC BY 4.0)}
\begin{document}
\maketitle

\noindent\emph{Source and licence.} Induced dynamics and quasifactors for minimal equicontinuous actions on Stone spaces. Version arXiv:2602.11034v1 (CC BY 4.0), distributed by arXiv under the Creative Commons Attribution 4.0 International licence, http://creativecommons.org/licenses/by/4.0/, as recorded in the arXiv licence metadata for this exact version and shown on the abstract page https://arxiv.org/abs/2602.11034v1. Full licence notice: \texttt{licenses/arxiv-2602.11034-CC-BY-4.0.txt}.\par\smallskip

\noindent\textit{Editorial scope.} Counted unit: the article's lemma lem:quasiFactorConstruction (source0.tex lines 987--999). The article's complete original proof of it is delivered with it (source0.tex lines 1000--1047, one \texttt{proof} environment of 48 lines). No mathematics is written, repaired or re-derived: the statement and the proof are the article's own lines, in the article's own notation, and no retained line is substituted or re-flowed.  No further source-local context is needed: every cross-reference of every retained line resolves inside the two delivered spans.  Nothing was omitted from any retained span, so the package carries no omission record.  No retained line contains a citation, so the wrapper carries no bibliography and no bibliography entry.  No retained line contains a figure environment, an image-inclusion command or drawing code, so no image is delivered and the package declares no asset dependency.  Editorial formatting only, in addition: an AMS \texttt{amsart} preamble that restates the article's own declarations and macros used by the retained lines, namely the environments \texttt{lemma} on separate counters, together with the standard packages those lines need.  The retained environments print in this document as Lemma 1 in order of appearance, whereas the article prints the same items with the numbering of its own full document.  The complete native source of the article as distributed in the arXiv e-print for this exact version is bundled unchanged as \texttt{sources/194456/source0.tex}.
\begin{lemma}
\label{lem:quasiFactorConstruction}
    Let $(X,G)$ be a minimal action, $g\in G$ and $\Gamma,\Lambda\in \Eig(X,G)$ such that $\Lambda\subseteq \Gamma$. 
    There exist factor maps $\pi_\Lambda\colon X\to G/\Lambda$ and 
    $\pi_\Gamma\colon X\to G/\Gamma$ with $\pi_\Gamma=\pi^\Gamma_\Lambda\circ \pi_\Lambda$. 
    Consider 
    \[A := \pi_\Gamma^{-1}(\Gamma)\cup \pi_\Lambda^{-1}(g\Lambda)\in \hyper(X).\]
\begin{itemize}
    \item[(i)] If $g\notin \Gamma$ and $[\Gamma\colon \Lambda]\geq 3$, then $G_0(A)=\Gamma\cap \Lambda^g$. 
    \item[(ii)] If $\Gamma$ is settled, $\Gamma^g\neq \Gamma$ and $[\Gamma\colon \Lambda]\geq 3$, then 
    $(\overline{G.A},G)$ is a finite quasifactor of $(X,G)$ that is not a factor.
\end{itemize}
\end{lemma}

\begin{proof}
    Since $\Lambda\in \Eig(X,G)$ there exists a factor map $\pi_\Lambda\colon X\to G/\Lambda$ and $\pi_\Gamma:=\pi^\Lambda_\Gamma\circ \pi_\Lambda$ yields a factor map $X\to G/\Gamma$. 

(i):
    Clearly, we have $G_0(A)\supseteq \Gamma\cap \Lambda^g$. 
    For the converse consider $h\in G_0(A)$. 
    Let $F\subseteq \Gamma$ be a fundamental domain of $\Lambda\leq \Gamma$ and denote
    $F':=F\cup \{g\}$.
    Note that
    \begin{align*}
        A
        =\pi^{-1}(\{g'\Lambda;\, g'\in F\}) \cup \pi^{-1}(g\Lambda)
        =\pi^{-1}\left(\{g'\Lambda;\, g'\in F'\}\right). 
    \end{align*}
    Thus, $h$ establishes a permutation $\sigma\colon F'\to F'$ with 
    $hg'\Lambda=\sigma(g')\Lambda$ for $g'\in F'$. 

    We next show that $h\in \Gamma$. 
    For this recall that $|F|=[\Gamma\colon \Lambda]\geq 3$.    
    Thus, there exist $g_1,g_2\in F$ with $\sigma(g_1)=g_2$. 
    We have $hg_1\in hg_1\Lambda = g_2\Lambda$ and hence 
    \[h\in g_2\Lambda g_1^{-1}\subseteq F\Lambda F\subseteq\Gamma\Gamma\Gamma=\Gamma.\] 

    It remains to show that $h\in \Lambda^g$. 
    For this recall that $g\in F'$. 
    If $\sigma(g)\neq g$ we have $\sigma(g)\in F$ and hence 
    \[g
    =h^{-1}hg\in h^{-1}hg\Lambda
    =h^{-1}\sigma(g)\Lambda
    \subseteq \Gamma F \Lambda
    \subseteq \Gamma\Gamma\Gamma=\Gamma,\]
    a contradiction. 
    This shows $\sigma(g)=g$ and it follows that
    \[h=hgg^{-1}\in hg\Lambda g^{-1}=\sigma(g)\Lambda g^{-1}=g\Lambda g^{-1}=\Lambda^g.\]

(ii):
    Note that $\Gamma^g\neq \Gamma$ implies $g\notin \Gamma$. 
    We thus observe from (i) that $G_0(A)=\Gamma\cap \Lambda^g$ is a finite index subgroup of $G$. 
    It follows that $(\overline{G.A},G)$ is finite. 

    If $(\overline{G.A},G)$ is a factor of $(X,G)$, then 
    $\Gamma\cap \Lambda^g=G_0(A)\in \Eig(X,G)$. 
    Since $\Eig(X,G)$ is upward closed we observe from 
    $\Gamma\cap \Lambda^g\subseteq \Gamma\cap \Gamma^g$ that 
    $\Gamma\cap \Gamma^g\in \Eig(X,G)$.
    From $\Gamma$ being settled it follows that $\Gamma=\Gamma^g$, a contradiction. 
    This shows that $(\overline{G.A},G)$ is not a factor of $(X,G)$. 
\end{proof}

\end{document}
